\documentclass[../../main.tex]{subfiles}

\begin{document}

\chapter{Routines}

% =====================================================
% =====================================================
% Rutina Raster Scan
% =====================================================
% =====================================================

\section{RasterScan} 

The routine performs a two-dimensional raster scan by moving along a selected axis while acquiring a profile along the perpendicular axis at each scan position. For example, when the X axis is selected as the scan axis, the system moves along X according to the specified scan step and acquires a Z-profile at each X position. This approach is particularly useful for vertical coupling measurements, where the optical power distribution can be characterized as a function of the X and Z positions. 

During the scan, the profiles acquired using the Suruga Seiki utilities can be saved for further analysis. The acquired data are then cleaned by removing invalid or irrelevant samples that could interfere with the generation of the heatmap. The processed profiles are interpolated onto a common grid and used to generate an interactive heatmap representing the measured power distribution. 

\begin{figure}[H]
    \centering
    \includegraphics[width=1\textwidth]{Routines/RasterScan/HowItWorks.png}
    \caption{Two-dimensional raster scan for optical coupling localization. Profiles are acquired along the profile axis while the stage is incrementally moved along the perpendicular scan axis. The acquired profiles are combined to reconstruct a two-dimensional heatmap, allowing the position of maximum optical coupling to be identified.}
    \label{fig:raster_scan}
\end{figure}

Once the scan is complete, the routine identifies the global maximum power measured within the scanned area and, if enabled, automatically moves the arm to the corresponding X and Z position.

\begin{figure}[H]

    \centering

    \includegraphics[width=1\textwidth]{Routines/RasterScan/BasicFlowchart.png}

    \caption{Flowchart of the \textit{RasterScan} routine. The process begins with the configuration of the raster scan and profile parameters, followed by the initialization of the Suruga Seiki platform. For each position along the scan axis, the profile and scan-axis positions are obtained, and the acquired data are retrieved, stored, and processed using \texttt{clean\_profile\_data()} and \texttt{save\_profile\_txt()}. The process is repeated until all positions along the scan axis have been completed. Finally, the processed profiles are passed to \texttt{create\_heatmap()} to reconstruct the two-dimensional power map and locate the maximum signal, after which the system can move to the corresponding target positions.}

    \label{fig:raster_scan_flowchart}

\end{figure}


% -----------------------------------------------------
% Descripcion funciones
% -----------------------------------------------------

% -----------------------------------------------------
\subsection{\texttt{clean\_profile\_data()}}
% -----------------------------------------------------

\textbf{Objective:}

Clean and prepare the data acquired by the profile function for further
processing.

This function receives the complete dictionary returned by
\texttt{profile.get\_profile\_data()} and extracts only the variables required
for the raster scan:

\begin{itemize}
    \item \texttt{main\_position}: the position of the axis being profiled.
    \item \texttt{signal1}: the measured optical signal.
\end{itemize}

The function then performs the following operations:

\begin{enumerate}
    \item Converts the extracted data into NumPy arrays.
    
    \item Removes invalid samples, including:
    \begin{itemize}
        \item Non-finite position values.
        \item Non-finite signal values.
        \item Samples where the position is equal to zero.
    \end{itemize}
    
    \item Sorts and reorganizes the remaining samples according to their
    position.
\end{enumerate}

The function \textbf{does not modify or save the original profile data}.
Its purpose is to generate a cleaner and more manageable representation of
the acquired data, which can then be used for further processing,
particularly for generating the heatmap with
\texttt{create\_heatmap()}.


% -----------------------------------------------------
\subsection{\texttt{save\_profile\_txt()}}
% -----------------------------------------------------

\textbf{Objective:}

Preserve the complete profile data acquired during the measurement, together
with the experimental parameters required to document and reproduce the
measurement.

This function receives the complete dictionary returned by
\texttt{profile.get\_profile\_data()}, including the following variables:

\begin{itemize}
    \item \texttt{main\_position}: position of the main axis during the
    profile.
    \item \texttt{sub1\_position}: position of the first secondary axis.
    \item \texttt{sub2\_position}: position of the second secondary axis.
    \item \texttt{signal1}: first measured signal.
    \item \texttt{signal2}: second measured signal.
\end{itemize}

In addition to the acquired data, the function constructs a header containing
the experimental parameters and the stage positions at the time of
acquisition:

\begin{itemize}
    \item \textbf{Date:} Date and time of the measurement.
    \item \textbf{Center X:} X-axis center position.
    \item \textbf{Center Z:} Z-axis center position.
    \item \textbf{Arm:} Selected measurement arm or configuration.
    \item \textbf{Scan axis:} Axis along which the raster scan is performed.
    \item \textbf{Scan width:} Total width of the scan.
    \item \textbf{Scan step:} Distance between consecutive scan positions.
    \item \textbf{Profile axis:} Axis along which the profile is acquired.
    \item \textbf{Profile range:} Total range covered by the profile.
    \item \textbf{Profile speed:} Speed used during profile acquisition.
\end{itemize}

The function then stores the complete dataset and the associated metadata in
a \textbf{tab-separated TXT file} using NumPy's
\texttt{np.savetxt()} function.

Unlike \texttt{clean\_profile\_data()}, this function does \textbf{not}
remove, filter, or reorganize the acquired samples. Its purpose is to
preserve the original measurement data as completely as possible, while
providing the experimental metadata necessary to identify and reproduce the
measurement.

The resulting TXT file therefore contains both the
\textbf{raw profile data} and the \textbf{measurement parameters} associated
with its acquisition.


% -----------------------------------------------------
\subsection{\texttt{create\_heatmap()}}
% -----------------------------------------------------

\textbf{Objective:}

Reconstruct a two-dimensional signal map from all the profiles acquired
during the raster scan, generate an interactive heatmap, and determine the
global maximum of \texttt{signal1} together with its corresponding X and Z
coordinates.

The function receives two inputs:

\begin{itemize}
    \item \texttt{scan\_data}: a dictionary containing the processed profiles
    acquired during the raster scan. Each entry has the following structure:
\end{itemize}

\begin{verbatim}
[
    {
        "x": ...,
        "z_positions": ...,
        "signal": ...
    },
    ...
]
\end{verbatim}

\begin{itemize}
    \item \texttt{output\_path}: the path where the generated heatmap will be
    saved.
\end{itemize}

\subsubsection{Data Reconstruction}

First, the function takes the Z positions from the first acquired profile
and uses them as the common reference axis:

\begin{verbatim}
z_common = scan_data[0]["z_positions"]
\end{verbatim}

Since the Z positions may not be identical for every profile, each profile is
interpolated onto this common Z-axis. This ensures that all profiles have the
same number of samples and can be represented consistently in a two-dimensional
matrix.

For each profile, the function:

\begin{enumerate}
    \item Retrieves the corresponding X position.
    \item Retrieves the Z positions and measured signal.
    \item Interpolates the signal onto \texttt{z\_common}.
    \item Adds the interpolated signal to the 2D signal matrix.
\end{enumerate}

The resulting matrix represents the measured signal as a function of the X
and Z coordinates and is used to generate the heatmap.


\subsubsection{Heatmap Generation}

The reconstructed 2D signal matrix is used to generate an
\textbf{interactive Plotly heatmap}, where:

\begin{itemize}
    \item The \textbf{X-axis} represents the raster scan position.
    \item The \textbf{Z-axis} represents the profile position.
    \item The \textbf{color scale} represents the measured
    \texttt{signal1}.
\end{itemize}

The resulting heatmap provides a two-dimensional visualization of the
measured optical signal across the scanned area and is saved to the location
specified by \texttt{output\_path}.


\subsubsection{Global Maximum Detection}

In addition to generating the heatmap, the function determines the position
of the maximum measured signal.

For each profile, the function first identifies its \textbf{local maximum}.
The local maximum is then compared with the maximum found in the previous
profiles. By continuously retaining the highest value, the function
ultimately determines the \textbf{global maximum} across the entire raster
scan.

The global maximum is stored together with its corresponding coordinates:

\begin{itemize}
    \item \textbf{Maximum signal:} maximum value of \texttt{signal1}.
    \item \textbf{X coordinate:} X position of the profile where the maximum
    was found.
    \item \textbf{Z coordinate:} Z position corresponding to the maximum
    signal within that profile.
\end{itemize}

Therefore, \texttt{create\_heatmap()} performs two main tasks: it
\textbf{reconstructs and visualizes the complete raster scan as a
two-dimensional signal map}, and it \textbf{identifies the position of the
maximum output power} within the scanned region.

\section{Automapper} 


% ===========================================
\subsection{checkPositionsAutomap}
% ===========================================

Validates the measurement plan generated by AutoMapper before moving the Suruga platform. It calculates and displays the absolute \texttt{x1} and \texttt{x2} positions for each device defined in the Excel file, but it does not connect to any instrument or move any axes.

\textbf{Source code:}
\href{https://github.com/niospinamev2/SurugaSeikiHex/blob/main/rutinas/checkPositionsAutomap.py}
{\texttt{checkPositionsAutomap.py}}.

\subsubsection{Workflow}

\begin{enumerate}
    \item Loads the chip layout and measurement plan from an Excel file using \texttt{automaper.load\_chip()}.

    \item Uses \texttt{origin\_x1} and \texttt{origin\_x2} as the reference origins.

    \item For each measurement, calculates the input and output offsets and the absolute positions of both arms.

    \item Displays the device, ports, and calculated positions, and waits for confirmation before displaying the next measurement.
\end{enumerate}

\subsubsection{Relevant Configuration}

%----------------------------------------------------------------
% Inicio tabla
%----------------------------------------------------------------

\footnotesize

\begin{longtblr}[
    caption = {Relevant configuration parameters for \texttt{checkPositionsAutomap}.},
    label = {tab:checkpositions-config}
]{
    colspec = {
        |Q[l, font=\ttfamily]|
        Q[l]|
        X[l]|
    },
    rowhead = 1,
    row{1} = {font=\bfseries},
    hlines,
}
%---------------------------
\SetRow{bg=gray9}
Parameter & Initial Value & Description \\
%---------------------------
\texttt{EXCEL\_FILE}
    & \texttt{chipMaps/SurugaTest.xlsx}
    & Chip layout and AutoMapper measurement plan. \\
%---------------------------
\texttt{origin\_x1} / \texttt{origin\_x2}
    & $0.0\,\mu\mathrm{m}$ / $0.0\,\mu\mathrm{m}$
    & Reference positions used to calculate the absolute coordinates.
\end{longtblr}

\normalsize

%----------------------------------------------------------------
% Fin tabla
%----------------------------------------------------------------

The absolute coordinates are calculated according to

\begin{equation}
    x_1 = \mathrm{origin}_{x1} - \mathrm{input}.x
\end{equation}

and

\begin{equation}
    x_2 = \mathrm{origin}_{x2} - \mathrm{output}.x.
\end{equation}

\subsubsection{Recommended Usage}

Run this routine whenever the Excel file, chip coupling, or origin convention changes. Compare the printed positions against the mechanical setup limits before running a routine that moves the machine.

\subsubsection{Results}

This routine does not create files or perform measurements. Its output consists of the position list displayed in the console.


% ===========================================
\subsection{moveAutomapOPM}
% ===========================================

Automatically iterates through the measurements defined by AutoMapper, positions the fibers at the input and output waveguides, performs alignment on both arms, and measures the average optical power using the optical power meter (OPM).

\textbf{Source code:}
\href{https://github.com/niospinamev2/SurugaSeikiHex/blob/main/rutinas/moveAutomapOPM.py}
{\texttt{moveAutomapOPM.py}}.

\subsection{Workflow}

\begin{enumerate}
    \item Loads the measurement plan from the AutoMapper Excel file.

    \item Connects to the Suruga platform and creates controllers for \texttt{x1}, \texttt{x2}, alignment, and the power meter.

    \item Loads \texttt{origin.txt} or uses the current position as the \texttt{WG0} origin.

    \item For each device, calculates the absolute positions of both arms and moves \texttt{x1} and \texttt{x2}.

    \item Performs a \textit{flat} alignment on the left arm followed by another alignment on the right arm.

    \item Reads ten samples from the OPM, calculates their average, and displays the result in the console.

    \item Requests confirmation before continuing and, once all measurements are completed, returns to \texttt{WG0}.
\end{enumerate}

\subsection{Relevant Configuration}

%----------------------------------------------------------------
% Inicio tabla
%----------------------------------------------------------------

\footnotesize

\begin{longtblr}[
    caption = {Relevant configuration parameters for \texttt{moveAutomapOPM}.},
    label = {tab:moveautomapopm-config}
]{
    colspec = {
        |Q[l, font=\ttfamily]|
        Q[l, font=\ttfamily]|
        X[l]|
    },
    rowhead = 1,
    row{1} = {font=\bfseries},
    hlines,
}

%---------------------------
\SetRow{bg=gray9}
Parameter & Initial Value & Description \\
%---------------------------
\texttt{EXCEL\_FILE}
    & \texttt{chipMaps/SurugaTest.xlsx}
    & Chip layout and measurement plan. \\
%---------------------------
\texttt{USE\_SAVED\_ORIGIN} / \texttt{SAVE\_ORIGIN}
    & \texttt{True} / \texttt{False}
    & Control whether \texttt{origin.txt} is loaded or saved. \\
%---------------------------
\texttt{ORIGIN\_FILE}
    & \texttt{origin.txt}
    & File containing the X1 and X2 coordinates of \texttt{WG0}. \\
%---------------------------
\texttt{flat\_left} / \texttt{flat\_right}
    & ---
    & Alignment parameters for each arm.
\end{longtblr}

\normalsize

%----------------------------------------------------------------
% Fin tabla
%----------------------------------------------------------------

\subsection{Results}

The average optical power for each device is printed to the console. The routine does not save the measurement values to a file.

\subsection{Precautions}

The routine moves both arms and requires manual confirmation between measurements. Verify the position plan using
\href{https://github.com/niospinamev2/SurugaSeikiHex/blob/main/rutinas/checkPositionsAutomap.py}{\texttt{checkPositionsAutomap}} before running it.


% ===========================================
\subsection{moveAutomapOPM\_OSA}
% ===========================================

Characterizes the devices defined by AutoMapper by combining Suruga platform
movement, automatic alignment, an average optical power measurement using the
OPM, and spectrum acquisition with the OSA.

\textbf{Source code:}
\href{https://github.com/niospinamev2/SurugaSeikiHex/blob/main/rutinas/moveAutomapOPM_OSA.py}
{\texttt{moveAutomapOPM\_OSA.py}}.

\subsection{Workflow}

\begin{enumerate}
    \item Loads the chip layout and measurement plan from Excel.

    \item Connects to the Suruga platform, OPM, and OSA, and configures the
    spectral sweep.

    \item Loads or saves the \texttt{WG0} origin in \texttt{origin2.txt}.

    \item For each device and repetition, positions \texttt{x1} and
    \texttt{x2} based on the corresponding input and output waveguides.

    \item Aligns the left and right arms.

    \item Runs an OSA sweep, retrieves the trace, and averages ten power
    readings.

    \item Saves the trace along with the device, ports, optical power, and
    repetition number as metadata.

    \item Returns to \texttt{WG0} and closes the OSA connection.
\end{enumerate}

\subsection{Relevant Configuration}

%----------------------------------------------------------------
% Inicio tabla
%----------------------------------------------------------------

\footnotesize

\begin{longtblr}[
    caption = {Relevant configuration parameters for
    \texttt{moveAutomapOPM\_OSA}.},
    label = {tab:moveautomapopm-osa-config}
]{
    colspec = {
        |Q[l, font=\ttfamily]|
        Q[l, font=\ttfamily]|
        X[l]|
    },
    rowhead = 1,
    row{1} = {font=\bfseries},
    hlines,
}

%---------------------------
\SetRow{bg=gray9}
Parameter & Initial Value & Description \\
%---------------------------
\texttt{DEBUG\_MODE}
    & \texttt{True}
    & Requests manual confirmations; when set to \texttt{False}, the routine
    runs automatically. \\
%---------------------------
\texttt{USE\_SAVED\_ORIGIN} / \texttt{SAVE\_ORIGIN}
    & \texttt{False} / \texttt{True}
    & Control whether the origin stored in \texttt{origin2.txt} is loaded or
    saved. \\
%---------------------------
\texttt{EXCEL\_FILE}
    & \texttt{chipMaps/SurugaTest01.xlsx}
    & Chip layout and AutoMapper measurement plan. \\
%---------------------------
\texttt{num\_repeticiones}
    & \texttt{1}
    & Number of complete passes through the measurement plan. \\
%---------------------------
\texttt{CENTER} / \texttt{SPAN}
    & 1550 nm / 90 nm
    & OSA spectral window. \\
%---------------------------
\texttt{RESOLUTION} / \texttt{SENSITIVITY}
    & 0.2 nm / -65 dBm
    & OSA resolution and sensitivity. \\
%---------------------------
\texttt{channel}
    & \texttt{1}
    & Power meter channel.
\end{longtblr}

\normalsize

%----------------------------------------------------------------
% Fin tabla
%----------------------------------------------------------------


\subsection{Results}

For each device, the routine saves a trace to:

\begin{center}
\texttt{data/<date>/medicion\_<repetition>\_<device>\_<input>\_<output>.txt}
\end{center}

The file contains the spectral data and metadata including the device, ports,
average optical power, and repetition number.

\subsection{Precautions}

With \texttt{DEBUG\_MODE=True}, the routine pauses before the OSA sweep,
between devices, and before returning to the origin.

Before using automatic mode, validate the movements with
\href{https://github.com/niospinamev2/SurugaSeikiHex/blob/main/rutinas/checkPositionsAutomap.py}
{\texttt{checkPositionsAutomap}}, verify the origin, and confirm that the OSA
VISA resource is available.

\end{document}


